% =====================================================================
% Section 9: discussion (key: disc). Written by the framing editor after all other sections.
% Sources: paper/CLAIMS.md "Answers to the question" and "Cross-track conflicts"; model notes-final §§9.2,
%   9.3, 10; universal notes-final §§11, 13; pa notes-final §§3.8, 7, 9.5; experiments notes-final §§11-13;
%   the section files (labels); the writers' reports of problems found while writing (see app-verification).
% No theorem environments. Literature: keys of bib/*.bib; unverified items hedged as in the sections.
% =====================================================================

\section{Discussion}\label{sec:disc}

The answers of \cref{sec:intro:answers} are conditional: the positive ones need data drawn from a law of the model, and the data the question has in mind, the theorems humans state, are not. This section says what the answers mean for the three refinements of the question, what a good version looks like and what it cannot promise, how the results relate to H\"anni's note and to the literature, and what remains open.

% ---------------------------------------------------------------------
\subsection{The three refinements}\label{sec:disc:proposal}

\paragraph{Templates.} H\"anni proposed templates to block his collapse. They block his schema: no non-ground $\DT$ template has all its instances in the acceptance part of the collapse schema (\cref{prop:time:notemplate}), and the reflection schema contains none either (\cref{prop:pa:refl}). They do not block the collapse at the level of theorems: in arithmetic one ground reflection sentence reproduces any $\Sigma_n$-sound assigner, at a cost linear in its length (\cref{prop:time:collapse}). Inside the class membership is matching, so likelihoods can be computed exactly over a pool (\cref{prop:exp:exact}). Templates have a second effect that the intuition behind the question does not anticipate. They bring the instance schema $\varphi(z)$ into the class, and $\varphi(z)$ is a simple law that implies exactly the observed instances; on closed instances $\forall x\varphi$ never overtakes it (\cref{thm:univ:B}). In a class of sentences only, under citation with the closure reading, $\forall x\varphi$ does take the posterior (\cref{prop:univ:sentences}(a), proof sketch). Nor do templates make the posterior see the logical boundaries of a schema: a schema and its root split tie exactly at matched weights, and with learned weights they are told apart by Occam terms or by usage statistics (\cref{prop:ident:splits,prop:ident:splitlzero,prop:ident:splitlone}).

\paragraph{A derivation-length prior.} The phrase has two readings. As a generative derivation grammar with a data-independent normaliser ($\Lone$), it supplies the size principle that H\"anni's 0/1 scores lack (\cref{prop:model:whichsize}); this is what makes memorisers and over-general templates lose. But $\Lone$'s code length counts grammar choices, not written symbols, and can be exponentially shorter than the written derivation (\cref{prop:time:lonesize}). If the length of a derivation is meant to price computation, it should be measured in written symbols ($\Lsig$, or H\"anni's graded score divided by its normaliser; \cref{rem:model:graded,prop:time:sigma}). The second reading, ``statements given without proof in our data set should be fairly easily derived from axioms'', makes the data theorems. Then a derivation likelihood favours the theory that makes the data cheap to state, not the logically minimal one. Under a two-part code this is memorisation for short theorems with long proofs (\cref{prop:pa:memo,rem:pa:streams}); the same pattern appears in the fragment $\{\mathrm{Q4},\mathrm{Q5}\}$ (\cref{prop:univ:eqfrag}, two-part comparison) and in the equational grammar $\Leq$ (\cref{ex:ident:lonedq}). The generator wins when the data come from its own derivation process (\cref{prop:pa:gibbs,prop:pa:wellspec}) or contain direct uses of its schemas. So a derivation-length prior identifies the working set of primitives of whoever produced the data, including frequently used lemmas (\cref{prop:univ:quant}(c), \cref{prop:pa:usage}), not a canonical axiomatisation.

\paragraph{A time penalty.} The penalty the question names, on the time ``to generate axioms / to check whether a statement is an axiom'', does not bite. Inside $\DT$ membership is matching (\cref{rem:model:priors}); and the collapse sets themselves have polynomial membership, so the brief's conjecture that such a penalty prices the collapse is refuted (\cref{prop:time:cheap}). What prices the collapse is derivation size in written symbols: every theory consistent with a hard assigner needs long derivations (\cref{thm:time:ntime,cor:time:hard}), and $\Lsig$ and the graded score charge this polynomially (\cref{prop:time:sigma}). The charge lies between the assigner's nondeterministic time, a proved lower bound, and its deterministic time, an upper bound with a proof sketch (\cref{rem:time:summary}); for plain $\Lone$ only a logarithmic charge is proved (\cref{cor:time:log,conj:time:poly}). Two other limits act differently. A bound on the time spent evaluating the likelihood sums over fewer derivations and so favours adopting theorems with long proofs as axioms (\cref{rem:time:links}(iv)). A steeper simplicity penalty trades early unsound lumping for faster removal of a false spare slot, and neither steepness removes both (\cref{rem:time:e7}).

% ---------------------------------------------------------------------
\subsection{A good version, and what it cannot promise}\label{sec:disc:good}

Collecting the pieces, a good version of the inducer has:
\begin{itemize}[leftmargin=*]
\item \emph{Hypotheses and prior}: finite sets of $\DT$ templates with a proper prefix-code prior and Dirichlet weights (\cref{def:model:prior,lem:model:kraft}); inconsistent theories removed by refutation at growing depth (\cref{prop:pa:refute}) or reported as a fourth value $\Inc$ (\cref{def:sound:tri}).
\item \emph{Likelihood}: a derivation grammar with a data-independent normaliser and a uniformly subcritical instantiation grammar (\cref{def:model:grammar,def:model:lone}), measured in written symbols ($\Lsig$) if it is to price computation; a model of which theorems get reported, by a known filter ($\Lsel$) or a prior over filters (\cref{thm:univ:omega}(e)); and a noise component or corruption channel, so that one false datum does not refute the true theory (\cref{prop:univ:robust}, \cref{rem:ident:nearmiss}).
\item \emph{Output}: the triple $(\Bel,\Dis,\Indep)$, a belief function with its plausibility, not a point estimate (\cref{sec:sound:tri}); and a verifier that accepts $s$ when $\Belc(s)\ge1-\delta$ (\cref{rem:sound:belc}), with a threshold that shrinks with $n$ when the weights are learned (\cref{thm:sound:shrink}).
\end{itemize}
Its likelihood is a computable real whose positivity is undecidable (\cref{prop:model:compute}); exact computation is feasible over finite pools with short derivations (\cref{sec:exp}).

\emph{What it can promise}, when the data are drawn from a law of the model: concentration on the generator class and convergence of predictions (\cref{thm:ident:doob,thm:ident:predded}); identification of the theorem set under $\Lone$ with parameters admissible, or of the instance union under $\Lzero$, with no bound on the number of templates and around Gold's theorem (\cref{cor:ident:deductive,thm:ident:limit,prop:ident:gold}); soundness of the verifier against adaptive provers with probability $1-\delta/\pi(C^*_d)$ for fixed weights, and the averaged or shrinking-threshold versions for learned weights (\cref{thm:sound:fixed,thm:sound:avg,thm:sound:shrink}).

\emph{What it cannot promise.} It cannot single out the actual axioms: inside the generator class the posterior is the prior (\cref{cor:ident:prior}). It cannot certify $\forall x\varphi$ from closed instances (\cref{thm:univ:omega}). It removes a false spare sentence only polynomially (\cref{prop:ident:spare}) and an inconsistent theory only by refutation (\cref{prop:pa:incons,prop:pa:refute}). Above all, it promises nothing when the data are not drawn from a law of the model, which is the case the question cares about: the posterior then goes to the best-fitting generator, weaker or unsound by turns (\cref{tab:ident:misspec}), and a verifier on it can be made to accept a falsehood with probability $1$ (\cref{prop:sound:misspec}). Under misspecification the soundness of the verifier depended, in the experiments, on what else was in the pool (\cref{rem:sound:e4}). The main lever is to model how the data were produced: which theorems are reported, how often each axiom form is used, and which mistakes occur. Each of these is a modelling choice that the data cannot check from positive instances alone.

% ---------------------------------------------------------------------
\subsection{H\"anni's note, point by point}\label{sec:disc:hanni}

His two variants are scores without the size principle, and ``does not contradict'' alone learns nothing beyond compatibility (\cref{prop:model:whichsize}(b), \cref{prop:pa:nc}). His weighting by provable givens and proof length, divided by its normaliser, is a two-part code that charges written derivation length; this makes good his remark that one can ``think of this as there being some model for how the statements get generated'' (\cref{rem:model:graded}). The trichotomy is best kept as a triple (\cref{sec:sound:tri}). The equivalence holds with the constant he expects, via Craig's trick, in a semimeasure convention (his note also considers speaking ``of a semimeasure here instead of renormalizing''); both weights depend on the data only as a set, so it holds for every reveal order (\cref{thm:time:equiv,def:time:inducers}). His ``easy'' separation from unrestricted function induction is \cref{prop:time:fiall}. Of his other ideas, ``only generating a very small subset of observations'' is what a likelihood with a data-independent normaliser does, with the selection modelled by $\Lsel$ or a prior over filters; ``making some mistakes'' is a noise mixture, harmless for soundness only if the noise model is part of the well-specified generator (E4(C1) in \cref{rem:sound:e4}); grading near misses is a corruption channel (\cref{rem:ident:nearmiss}). His time-bounded note has an analogue in a derivation-bounded inducer (\cref{conj:time:polytime}).

% ---------------------------------------------------------------------
\subsection{Relation to the literature}\label{sec:disc:lit}

We cite and do not re-prove; \cref{app:ver:refs} lists the references that were not checked against their sources. \Cref{thm:univ:confirm} is the countable-class analogue of the confirmation of universal hypotheses by the universal prior \citep{hutter2007universal}; \cref{thm:univ:omega} shows that what gets confirmed is ``all future data are instances'', not derivability. Confirmation is not stepwise monotone \citep{leike2015solomonoff}, and for truth rather than derivability Gaifman's condition applies (\cref{prop:univ:gaifman}; \citealp{gaifman1964concerning}). \Cref{thm:ident:doob} is the argument of \citet{doob1949application} for a countable class; consistency in larger models goes back to \citet{schwartz1965bayes} (see \citealp{ghosal2017fundamentals}), and misspecified posteriors concentrate on KL-minimisers in general \citep{berk1966limiting,kleijn2006misspecification}, under conditions we did not check for template classes. Learning from text is limited by Gold's theorem \citep{gold1967language,angluin1980inductive}; learning from stochastic positive data goes back to \citet{horning1969study} and \citet{angluin1988identifying}. The size principle is that of \citet{tenenbaum2001generalization}, the code-length reading is MDL \citep{rissanen1978modeling} with the KT estimator \citep{krichevsky1981performance}, and the spare-slot rates are of the kind studied for overfitted mixtures \citep{rousseau2011asymptotic}. $\Bel$ and $\Pl$ are a belief function and its plausibility in the sense of \citet{shafer1976mathematical}. The collapse results use Craig's trick \citep{craig1953axiomatizability} and the nondeterministic time hierarchy \citep{cook1972hierarchy,seiferas1978separating,zak1983turing}; Levin's Kt and the speed prior enter only as examples of membership- or generation-time penalties; the upper bounds recall \citet{pudlak1998lengths}; Kreisel's conjecture is related (\cref{rem:univ:kreisel}). Inductive logic programming \citep{muggleton1991inductive} learns clauses from examples and background knowledge; AS is the non-Bayesian counterpart of this paper, and IL supplies its Ville argument \citep{claude2026whatfollows,ville1939etude}.

% ---------------------------------------------------------------------
\subsection{Open problems}\label{sec:disc:open}

Collected from the four tracks and from the writing of this paper, and deduplicated; the cited statements give the details.
\begin{enumerate}[leftmargin=*,itemsep=1pt]
\item \emph{Rates}: identification for every weight vector (\cref{thm:ident:limit}(c)); rigorous spare-slot rates (\cref{prop:ident:spare}); the Occam rate under $\Lone$ (\cref{prop:ident:splitlone}(d)); the memoriser rate (\cref{prop:univ:memo}(iii)).
\item \emph{Competitors only $\Lone$ keeps alive}: partial splits that derive the rest (\cref{rem:ident:mdl}); whether the two-part maximum can change a deductive limit (\cref{prop:model:max}); whether $\Lone$ separates equivalent theories for all grammar parameters (\cref{prop:ident:sep}).
\item \emph{Misspecification}: unsound mass in the full class (\cref{rem:ident:fullclass}); the KL-minimising families for instance-only data; \cref{conj:ident:cder}; soundness bounds in terms of the distance from the data law to the pool, and robust posteriors.
\item \emph{The verifier}: completeness under the shrinking threshold; strictness of the inclusion in \cref{prop:sound:cautious}(d).
\item \emph{$\forall x\varphi$}: reversal through $\wedge$-detours (\cref{rem:univ:detour}); the race between split and memoriser families over $\Q$ (\cref{prop:univ:noguard,prop:univ:sentences}); an unknown selection learned jointly with the templates (\cref{thm:univ:omega}(e)); the full-sum comparison at the default law (\cref{prop:univ:eqfrag}); formulas with several variables.
\item \emph{The collapse}: \cref{conj:time:poly}; whether axiom induction with a symbol-size penalty equals function induction over assigners with short certificates up to polynomial factors; a single-sorted repair of H\"anni's schema; the $\ISigma n$ comparison under a derivation-time penalty; \cref{conj:time:polytime}.
\item \emph{$\PA$ and $\ZF$}: the posterior over all template unions; non-read-once induction templates (\cref{conj:pa:broad}); theorem data under plain $\Lone$, the full sum or Dirichlet weights; Collection from Replacement without Foundation.
\item \emph{Computation}: search beyond a pool; tree derivations in code, so that $\PA$ can be run under $\Lone$; a learned instantiation grammar; cheap remedies for early unsound lumps (\cref{rem:exp:limits}).
\item \emph{Near misses}: how much natural corruption channels enlarge the generator class (\cref{rem:ident:nearmiss}).
\item \emph{Write-ups}: full proofs of the proof sketches listed in \cref{app:ver:sketches}.
\end{enumerate}

\paragraph{Closing.} This was a check of an idea, not the design of a prover. The idea survives in a precise form: a template prior with a generative derivation likelihood is a coherent inducer, it confirms predictions, and with well-specified data it finds a system with the right theorems (or instances) and supports a verifier that adaptive provers fool only with small probability. What it does not do is turn human-stated theorems into the axioms behind them. The gap lies between how the model generates data and how mathematicians select what they write down.
